\section{Related Work}
\label{sec:related}

\paragraph{State legislative redistricting literature.}
The political science literature on state legislative redistricting has focused
primarily on detecting and measuring partisan gerrymandering.
\citet{chen2013unintentional} demonstrated using simulation that geographic
sorting of Democratic voters into dense urban areas produces systematic
partisan bias in state legislative maps even without intentional gerrymandering.
\citet{stephanopoulos2015partisan} developed the efficiency gap as a measure
of partisan gerrymandering and applied it to state legislative maps, finding
that the 2010 round of redistricting produced the most extreme efficiency gaps
in modern American history in states such as Wisconsin, North Carolina, and Ohio.
\citet{mcghee2014measuring} earlier characterised partisan bias in
single-member systems using seat-vote curves, noting that state legislative
systems are more sensitive to partisan bias than congressional systems due to
the larger number of smaller districts.

\paragraph{Algorithmic approaches.}
\citet{cirincione2000assessing} and \citet{altman2011bard} pioneered
computer-generated redistricting for state legislative applications,
using simulated annealing to generate random baseline maps.
\citet{fifield2020automated} developed the sequential Monte Carlo (SMC)
algorithm for state legislative redistricting, enabling exploration of
the plan space for chambers with up to $k \approx 100$ districts.
Our approach differs from these in its objective: rather than exploring
the plan space, we seek the single best plan under the minimum-edge-cut
criterion, making the algorithmic output directly comparable to a
proposed enacted plan.

\paragraph{Compactness in state legislative redistricting.}
Many state constitutions include explicit compactness requirements for
legislative maps \citep{monmonier2001bushmanders}. The Polsby-Popper
measure \citep{polsby1991third} and the Reock score \citep{reock1961note}
are the most commonly used metrics in legal and academic contexts.
\citet{kaufman2020state} surveys state compactness standards and finds
wide variation in both the stringency and the operationalisation of
compactness requirements. Our work provides a uniform algorithmic baseline
for comparison across all states.

\paragraph{Block-group resolution.}
Prior work on block-group redistricting is sparse. \citet{browdy1990computer}
noted that block-level data enables finer population balance but at high
computational cost. The MAUP implications of resolution choice for
redistricting outcomes are studied in \citet{openshaw1984modifiable}
and, in the congressional context, in companion paper C.1
\citep{dellaLibera2026maup}. We extend the resolution analysis to
state legislative redistricting in companion paper F.3.
